\begin{table}[h!]
	\centering
	\caption{Institutional Manifestations of External Constraint Regimes (Chile, 1943--2010)}
	\label{tab:chile_binding_regimes}
	\small
	\begin{tabular}{p{2.0cm} p{1.8cm} p{1.0cm} p{7.2cm}}
		\hline
		Historical Period & Regime Status ($D_t^\gamma$) & Duration & Institutional Manifestation of the External Constraint \\
		\hline
		1943--1984 & Binding ($D_t = 1$) & 42 years & ISI Reserve Scarcity \& Structural Inflation: Established under Law 7,200 foreign-exchange budgeting and fixed/crawling pegs; reserve depletion ($S_t \to 0$) under wage pressure forced maxidevaluations, cost-plus markups, and structural inflation spirals. \\
		\hline
		1985--2005 & Slack ($D_t = 0$) & 21 years & Export Consolidation \& Reserve Buffering: Post-1982 debt-crisis and recovery non-traditional export expansion, Concertación democratic transition, large FDI inflows, and central bank foreign reserve accumulation. \\
		\hline
		2006--2010 & Binding ($D_t = 1$) & 5 years & Commodity Boom, Rentier Surplus Drain \& Re-primarization: Under a free-floating exchange rate and inflation targeting (3\%), FX starvation did not trigger inflation; instead, high import demand ($\tilde{M}_t$) and massive MNC profit repatriation (DL 600) manifested as a rentist surplus drain and Dutch Disease re-primarization. \\
		\hline
		\multicolumn{4}{p{14.2cm}}{\footnotesize \textit{Notes:} Classification based on the optimal profile threshold $\hat{\gamma} = 0.425$ evaluated over the standardized composite balance-of-payments index $Z_t = 0.2 \tilde{M}_t + 0.8 \tilde{S}_t$. Regime $D_t = 1$ denotes a binding balance-of-payments constraint where foreign-exchange starvation forces peripheral firms into a corner solution ($\bar{\kappa}_t$). Total sample size $T = 68$ (1943--2010 effective cointegrating window), yielding 47 binding years ($69.1\%$) and 21 slack years ($30.9\%$).} \\
	\end{tabular}
\end{table}